% =====================================================================
% Abbildung: Modellkette: Von der Anforderung zum Testfall (Hinderniserkennung)
% Passend zu Kapitel 2, Abschnitt "Beispielkette: Hinderniserkennung"
% (Unterabschnitt von "MBSE im Roboterprojekt")
%
% Einbindung im Buch, z. B. in chapters/02_mbse.tex:
%
%   \begin{figure}[H]
%     \centering
%     % =====================================================================
% Abbildung: Modellkette: Von der Anforderung zum Testfall (Hinderniserkennung)
% Passend zu Kapitel 2, Abschnitt "Beispielkette: Hinderniserkennung"
% (Unterabschnitt von "MBSE im Roboterprojekt")
%
% Einbindung im Buch, z. B. in chapters/02_mbse.tex:
%
%   \begin{figure}[H]
%     \centering
%     % =====================================================================
% Abbildung: Modellkette: Von der Anforderung zum Testfall (Hinderniserkennung)
% Passend zu Kapitel 2, Abschnitt "Beispielkette: Hinderniserkennung"
% (Unterabschnitt von "MBSE im Roboterprojekt")
%
% Einbindung im Buch, z. B. in chapters/02_mbse.tex:
%
%   \begin{figure}[H]
%     \centering
%     % =====================================================================
% Abbildung: Modellkette: Von der Anforderung zum Testfall (Hinderniserkennung)
% Passend zu Kapitel 2, Abschnitt "Beispielkette: Hinderniserkennung"
% (Unterabschnitt von "MBSE im Roboterprojekt")
%
% Einbindung im Buch, z. B. in chapters/02_mbse.tex:
%
%   \begin{figure}[H]
%     \centering
%     \input{figures/fig_kap02_modellkette_traceability}
%     \caption{Farbkodierte Modellkette von einem Stakeholderbedarf ueber
%       Anforderung, Funktionen und Systemelemente bis zu ROS2-Node und
%       Testfall.}
%     \label{fig:kap02-modellkette-traceability}
%   \end{figure}
%
% Benoetigte Pakete (im Buch bereits in main.tex geladen): tikz
% Benoetigte TikZ-Bibliotheken: positioning, arrows.meta, shadows, calc
% =====================================================================

\input{figures/fig_kap02_modellkette_traceability.tikz}

%     \caption{Farbkodierte Modellkette von einem Stakeholderbedarf ueber
%       Anforderung, Funktionen und Systemelemente bis zu ROS2-Node und
%       Testfall.}
%     \label{fig:kap02-modellkette-traceability}
%   \end{figure}
%
% Benoetigte Pakete (im Buch bereits in main.tex geladen): tikz
% Benoetigte TikZ-Bibliotheken: positioning, arrows.meta, shadows, calc
% =====================================================================

\begin{tikzpicture}[
    every node/.style={font=\small},
    ext/.style={
        rectangle, rounded corners=1.5pt, draw=black!50, thick, dashed,
        minimum width=24mm, minimum height=13mm,
        align=center, fill=white, drop shadow, font=\scriptsize
    },
    req/.style={
        rectangle, rounded corners=1.5pt, draw=black!70, thick,
        minimum width=24mm, minimum height=13mm,
        align=center, fill=blue!12, drop shadow, font=\scriptsize
    },
    part/.style={
        rectangle, rounded corners=1.5pt, draw=green!40!black, thick,
        minimum width=24mm, minimum height=13mm,
        align=center, fill=green!15, drop shadow, font=\scriptsize
    },
    tool/.style={
        rectangle, rounded corners=1.5pt, draw=orange!60!black, thick,
        minimum width=24mm, minimum height=13mm,
        align=center, fill=orange!15, drop shadow, font=\scriptsize
    },
    test/.style={
        rectangle, rounded corners=1.5pt, draw=black!50, thick,
        minimum width=24mm, minimum height=13mm,
        align=center, fill=black!6, drop shadow, font=\scriptsize
    },
    flow/.style={-{Stealth[length=2.2mm]}, thick, black!60}
]

% --- Zeile 1: Stakeholderbedarf bis Systemelement Lidar-Sensor ---
\node[ext]  (need) at (1.2,3.5)  {Stakeholder-\\bedarf};
\node[req]  (req)  at (4.0,3.5)  {Anforderung};
\node[part] (f1)   at (6.8,3.5)  {Funktion:\\Hindernis\\erkennen};
\node[part] (f2)   at (9.6,3.5)  {Funktion:\\Bewegung sicher\\begrenzen};
\node[part] (lidar)at (12.4,3.5) {Systemelement:\\Lidar-Sensor};

\draw[flow] (need) -- (req);
\draw[flow] (req)  -- (f1);
\draw[flow] (f1)   -- (f2);
\draw[flow] (f2)   -- (lidar);

% --- Zeile 2: Systemelement Sicherheitslogik bis Testfall ---
\node[part] (safety) at (12.4,0) {Systemelement:\\Sicherheitslogik};
\node[tool] (topic)  at (9.6,0)  {ROS2-Topic};
\node[tool] (rosnode)   at (6.8,0)  {ROS2-Node};
\node[test] (tc)     at (4.0,0)  {Testfall};

\draw[flow] (lidar.south)  -- (safety.north);
\draw[flow] (safety) -- (topic);
\draw[flow] (topic)  -- (rosnode);
\draw[flow] (rosnode)   -- (tc);

% --- Legende ---
\node[req, minimum width=6mm, minimum height=4mm, font=\tiny] (leg1) at (1.2,-2.2) {};
\node[right=1mm of leg1, font=\tiny, align=left] {Anforderung};
\node[part, minimum width=6mm, minimum height=4mm, font=\tiny, right=18mm of leg1] (leg2) {};
\node[right=1mm of leg2, font=\tiny, align=left] {Funktion / Systemelement};
\node[tool, minimum width=6mm, minimum height=4mm, font=\tiny, right=50mm of leg1] (leg3) {};
\node[right=1mm of leg3, font=\tiny, align=left] {ROS2-Element};
\node[test, minimum width=6mm, minimum height=4mm, font=\tiny, right=7.48cm of leg1] (leg4) {};
\node[right=1mm of leg4, font=\tiny, align=left] {Testfall};

\end{tikzpicture}


%     \caption{Farbkodierte Modellkette von einem Stakeholderbedarf ueber
%       Anforderung, Funktionen und Systemelemente bis zu ROS2-Node und
%       Testfall.}
%     \label{fig:kap02-modellkette-traceability}
%   \end{figure}
%
% Benoetigte Pakete (im Buch bereits in main.tex geladen): tikz
% Benoetigte TikZ-Bibliotheken: positioning, arrows.meta, shadows, calc
% =====================================================================

\begin{tikzpicture}[
    every node/.style={font=\small},
    ext/.style={
        rectangle, rounded corners=1.5pt, draw=black!50, thick, dashed,
        minimum width=24mm, minimum height=13mm,
        align=center, fill=white, drop shadow, font=\scriptsize
    },
    req/.style={
        rectangle, rounded corners=1.5pt, draw=black!70, thick,
        minimum width=24mm, minimum height=13mm,
        align=center, fill=blue!12, drop shadow, font=\scriptsize
    },
    part/.style={
        rectangle, rounded corners=1.5pt, draw=green!40!black, thick,
        minimum width=24mm, minimum height=13mm,
        align=center, fill=green!15, drop shadow, font=\scriptsize
    },
    tool/.style={
        rectangle, rounded corners=1.5pt, draw=orange!60!black, thick,
        minimum width=24mm, minimum height=13mm,
        align=center, fill=orange!15, drop shadow, font=\scriptsize
    },
    test/.style={
        rectangle, rounded corners=1.5pt, draw=black!50, thick,
        minimum width=24mm, minimum height=13mm,
        align=center, fill=black!6, drop shadow, font=\scriptsize
    },
    flow/.style={-{Stealth[length=2.2mm]}, thick, black!60}
]

% --- Zeile 1: Stakeholderbedarf bis Systemelement Lidar-Sensor ---
\node[ext]  (need) at (1.2,3.5)  {Stakeholder-\\bedarf};
\node[req]  (req)  at (4.0,3.5)  {Anforderung};
\node[part] (f1)   at (6.8,3.5)  {Funktion:\\Hindernis\\erkennen};
\node[part] (f2)   at (9.6,3.5)  {Funktion:\\Bewegung sicher\\begrenzen};
\node[part] (lidar)at (12.4,3.5) {Systemelement:\\Lidar-Sensor};

\draw[flow] (need) -- (req);
\draw[flow] (req)  -- (f1);
\draw[flow] (f1)   -- (f2);
\draw[flow] (f2)   -- (lidar);

% --- Zeile 2: Systemelement Sicherheitslogik bis Testfall ---
\node[part] (safety) at (12.4,0) {Systemelement:\\Sicherheitslogik};
\node[tool] (topic)  at (9.6,0)  {ROS2-Topic};
\node[tool] (rosnode)   at (6.8,0)  {ROS2-Node};
\node[test] (tc)     at (4.0,0)  {Testfall};

\draw[flow] (lidar.south)  -- (safety.north);
\draw[flow] (safety) -- (topic);
\draw[flow] (topic)  -- (rosnode);
\draw[flow] (rosnode)   -- (tc);

% --- Legende ---
\node[req, minimum width=6mm, minimum height=4mm, font=\tiny] (leg1) at (1.2,-2.2) {};
\node[right=1mm of leg1, font=\tiny, align=left] {Anforderung};
\node[part, minimum width=6mm, minimum height=4mm, font=\tiny, right=18mm of leg1] (leg2) {};
\node[right=1mm of leg2, font=\tiny, align=left] {Funktion / Systemelement};
\node[tool, minimum width=6mm, minimum height=4mm, font=\tiny, right=50mm of leg1] (leg3) {};
\node[right=1mm of leg3, font=\tiny, align=left] {ROS2-Element};
\node[test, minimum width=6mm, minimum height=4mm, font=\tiny, right=7.48cm of leg1] (leg4) {};
\node[right=1mm of leg4, font=\tiny, align=left] {Testfall};

\end{tikzpicture}


%     \caption{Farbkodierte Modellkette von einem Stakeholderbedarf ueber
%       Anforderung, Funktionen und Systemelemente bis zu ROS2-Node und
%       Testfall.}
%     \label{fig:kap02-modellkette-traceability}
%   \end{figure}
%
% Benoetigte Pakete (im Buch bereits in main.tex geladen): tikz
% Benoetigte TikZ-Bibliotheken: positioning, arrows.meta, shadows, calc
% =====================================================================

\begin{tikzpicture}[
    every node/.style={font=\small},
    ext/.style={
        rectangle, rounded corners=1.5pt, draw=black!50, thick, dashed,
        minimum width=24mm, minimum height=13mm,
        align=center, fill=white, drop shadow, font=\scriptsize
    },
    req/.style={
        rectangle, rounded corners=1.5pt, draw=black!70, thick,
        minimum width=24mm, minimum height=13mm,
        align=center, fill=blue!12, drop shadow, font=\scriptsize
    },
    part/.style={
        rectangle, rounded corners=1.5pt, draw=green!40!black, thick,
        minimum width=24mm, minimum height=13mm,
        align=center, fill=green!15, drop shadow, font=\scriptsize
    },
    tool/.style={
        rectangle, rounded corners=1.5pt, draw=orange!60!black, thick,
        minimum width=24mm, minimum height=13mm,
        align=center, fill=orange!15, drop shadow, font=\scriptsize
    },
    test/.style={
        rectangle, rounded corners=1.5pt, draw=black!50, thick,
        minimum width=24mm, minimum height=13mm,
        align=center, fill=black!6, drop shadow, font=\scriptsize
    },
    flow/.style={-{Stealth[length=2.2mm]}, thick, black!60}
]

% --- Zeile 1: Stakeholderbedarf bis Systemelement Lidar-Sensor ---
\node[ext]  (need) at (1.2,3.5)  {Stakeholder-\\bedarf};
\node[req]  (req)  at (4.0,3.5)  {Anforderung};
\node[part] (f1)   at (6.8,3.5)  {Funktion:\\Hindernis\\erkennen};
\node[part] (f2)   at (9.6,3.5)  {Funktion:\\Bewegung sicher\\begrenzen};
\node[part] (lidar)at (12.4,3.5) {Systemelement:\\Lidar-Sensor};

\draw[flow] (need) -- (req);
\draw[flow] (req)  -- (f1);
\draw[flow] (f1)   -- (f2);
\draw[flow] (f2)   -- (lidar);

% --- Zeile 2: Systemelement Sicherheitslogik bis Testfall ---
\node[part] (safety) at (12.4,0) {Systemelement:\\Sicherheitslogik};
\node[tool] (topic)  at (9.6,0)  {ROS2-Topic};
\node[tool] (rosnode)   at (6.8,0)  {ROS2-Node};
\node[test] (tc)     at (4.0,0)  {Testfall};

\draw[flow] (lidar.south)  -- (safety.north);
\draw[flow] (safety) -- (topic);
\draw[flow] (topic)  -- (rosnode);
\draw[flow] (rosnode)   -- (tc);

% --- Legende ---
\node[req, minimum width=6mm, minimum height=4mm, font=\tiny] (leg1) at (1.2,-2.2) {};
\node[right=1mm of leg1, font=\tiny, align=left] {Anforderung};
\node[part, minimum width=6mm, minimum height=4mm, font=\tiny, right=18mm of leg1] (leg2) {};
\node[right=1mm of leg2, font=\tiny, align=left] {Funktion / Systemelement};
\node[tool, minimum width=6mm, minimum height=4mm, font=\tiny, right=50mm of leg1] (leg3) {};
\node[right=1mm of leg3, font=\tiny, align=left] {ROS2-Element};
\node[test, minimum width=6mm, minimum height=4mm, font=\tiny, right=7.48cm of leg1] (leg4) {};
\node[right=1mm of leg4, font=\tiny, align=left] {Testfall};

\end{tikzpicture}

